\documentclass[11pt]{article}

\usepackage{amsmath}
\usepackage{geometry}
\geometry{margin=1in}

\title{AI Opportunity Atlas v1\\Methodology}
\author{}
\date{}

\begin{document}

\maketitle

\section{Core Idea}

Each scientific problem is treated as an ``opportunity node'' characterized by:

\begin{itemize}
    \item Impact (human and economic value)
    \item Difficulty (scientific or engineering resistance)
    \item AI solvability (likelihood that frontier AI systems make progress)
    \item Leverage (how much solving the problem accelerates other science)
\end{itemize}

The goal is to rank problems by expected value of AI-enabled scientific acceleration.

\section{Primary Variables}

\subsection{Humanity Impact (H)}
\textbf{Scale:} 0--100

Represents:
\begin{itemize}
    \item Lives saved or improved
    \item Scientific acceleration potential
    \item Societal transformation potential
\end{itemize}

Examples:
AGI alignment $\approx 100$, protein design $\approx 99$

\subsection{Economic Impact (E)}
\textbf{Scale:} 0--100

Represents:
\begin{itemize}
    \item GDP-level economic impact
    \item Industrial transformation
    \item Cost reduction potential
\end{itemize}

\subsection{Scientific Difficulty (D)}
\textbf{Scale:} 0--100

Represents:
\begin{itemize}
    \item Theoretical hardness
    \item Experimental cost
    \item Computational irreducibility
    \item Time-to-breakthrough complexity
\end{itemize}

Examples:
P vs NP $\approx 100$, protein design $\approx 70$--85

\subsection{LLM Progress Likelihood (L)}
\textbf{Scale:} 0--100

Represents probability that frontier LLMs or multimodal AI systems can make meaningful progress within 5--10 years.

\begin{itemize}
    \item High: protein design, interpretability
    \item Low: pure math conjectures, deep physics proofs
\end{itemize}

\section{Primary Scoring Function}

The main ranking metric is:

\begin{equation}
\text{AILeverageScore} = \frac{H \times L}{D}
\end{equation}

Where:
\begin{itemize}
    \item $H$ = Humanity Impact
    \item $L$ = LLM Progress Likelihood
    \item $D$ = Scientific Difficulty
\end{itemize}

\section{Interpretation}

This score measures:

\begin{quote}
How much societal value can plausibly be unlocked per unit of scientific difficulty using AI systems.
\end{quote}

\section{Secondary Derived Metrics}

\subsection{Expected Societal Return (ESR)}

\begin{equation}
ESR = H \times E
\end{equation}

Captures total potential societal and economic value.

\subsection{AI Replaceability Index (ARI)}

A heuristic measure of how much a problem is:

\begin{itemize}
    \item search-based
    \item pattern-based
    \item simulation-based
\end{itemize}

Higher values indicate better suitability for AI-driven solutions.

\subsection{Recursive AI Benefit (RAB)}

Measures whether solving the problem improves AI systems themselves.

Examples:
\begin{itemize}
    \item Interpretability: high
    \item Alignment: high
    \item Theorem proving: high
    \item Climate modeling: low-medium
\end{itemize}

\section{Problem Taxonomy}

Each entry is classified as:

\begin{itemize}
    \item OP: Explicit open problem or conjecture
    \item TECH: Foundational technical paper revealing bottleneck
    \item REV: Survey of unresolved challenges
    \item ROAD: Research roadmap or agenda
\end{itemize}

\section{Final Ranking Rule}

Final ranking is determined by a multi-objective heuristic:

\begin{equation}
Rank \approx f(\text{AILeverageScore}, ESR, RAB, \text{confidence})
\end{equation}

Where:
\begin{itemize}
    \item AILeverageScore dominates primary ordering
    \item ESR biases toward high economic impact
    \item RAB increases priority for AI-self-improving problems
    \item Confidence adjusts for epistemic uncertainty
\end{itemize}

\section{Key Limitation}

This framework is a heuristic model, not a learned or calibrated predictive system.

Limitations include:
\begin{itemize}
    \item Subjective scoring estimates
    \item No uncertainty distributions
    \item No empirical calibration to historical breakthroughs
    \item Correlated variable bias
\end{itemize}

\section{Summary}

The AI Opportunity Atlas ranks scientific problems by expected value of AI-enabled search acceleration, normalized by intrinsic difficulty, with additional emphasis on recursive AI improvement potential.

\end{document}